\begin{question}
\noindent $A$, $B$, $C$ and $D$ are four points on the circumference of a circle, forming a cyclic quadrilateral $ABCD$.

\noindent Angle $A = 2x$ and angle $C = 3x + 10$, as shown in the diagram below.

\begin{center}
\begin{tikzpicture}[scale=2]
    \coordinate[label=left:$A$] (A) at (150:1);
    \coordinate[label=above right:$B$] (B) at (60:1);
    \coordinate[label=right:$C$] (C) at (-30:1);
    \coordinate[label=below left:$D$] (D) at (220:1);

    \draw[name path=circ] (0,0) circle (1);
    \draw[line join=bevel] (A) -- (B) -- (C) -- (D) -- cycle;

    \pic [draw, "$2x$", angle eccentricity=1.5, angle radius = 0.5cm] {angle = D--A--B};
    \pic [draw, "$3x+10$", angle eccentricity=1.4, angle radius = 0.5cm] {angle = B--C--D};
\end{tikzpicture}
\end{center}

\begin{enumerate}
    \item[(a)] Calculate the value of $x$.

    \vspace{4cm}
    \answereq{x}{2}

    \item[(b)] Calculate the size of angle $A$ and the size of angle $C$.

    \vspace{3cm}
    \answerlines{2}{2}
\end{enumerate}
\end{question}
